% Title: Explorations of the Cauchy Momentum Equation.
% Author: Amlal El Mahrouss (A. E. Mahrouss)

\documentclass[10pt]{article}

%% ---- Core Packages ----------------------------------------
\usepackage[T1]{fontenc}
\usepackage[utf8]{inputenc}
\usepackage{lmodern}
\usepackage{microtype}

%% ---- Page Geometry ----------------------------------------
\usepackage[
  top=1in, bottom=1in,
  left=0.75in, right=0.75in,
  columnsep=0.25in
]{geometry}

%% ---- Mathematics ------------------------------------------
\usepackage{amsmath, amssymb, amsthm, amsfonts}
\usepackage{mathtools}
\usepackage{bm}
\usepackage{textcomp}
\usepackage{esint}

%% ---- Figures & Tables -------------------------------------
\usepackage{graphicx}
\usepackage{booktabs}
\usepackage{multirow}
\usepackage{array}
\usepackage{caption}
\usepackage{subcaption}
\usepackage{float}

%% ---- Code Listings ----------------------------------------
\usepackage{listings}
\usepackage{xcolor}
\usepackage{algorithm}
\usepackage{algpseudocode}

\lstdefinestyle{codestyle}{
  backgroundcolor=\color{gray!8},
  basicstyle=\ttfamily\footnotesize,
  breakatwhitespace=false,
  breaklines=true,
  captionpos=b,
  commentstyle=\color{green!50!black},
  keywordstyle=\color{blue}\bfseries,
  stringstyle=\color{orange!70!black},
  numberstyle=\tiny\color{gray},
  numbers=left,
  numbersep=5pt,
  showstringspaces=false,
  frame=single,
  rulecolor=\color{gray!40},
  tabsize=2
}
\lstset{style=codestyle}

%% ---- Hyperlinks -------------------------------------------
\usepackage[
  colorlinks=true,
  linkcolor=blue!70!black,
  citecolor=green!50!black,
  urlcolor=blue!60!black
]{hyperref}
\usepackage{url}

%% ---- Bibliography -----------------------------------------
\usepackage[numbers, sort&compress]{natbib}

%% ---- Miscellaneous ----------------------------------------
\usepackage{enumitem}
\usepackage{siunitx}
\usepackage{cleveref}

%% ---- Theorem Environments ---------------------------------
\theoremstyle{plain}
\newtheorem{theorem}{Theorem}[section]
\newtheorem{lemma}[theorem]{Lemma}
\newtheorem{corollary}[theorem]{Corollary}
\newtheorem{proposition}[theorem]{Proposition}

\theoremstyle{definition}
\newtheorem{definition}[theorem]{Definition}
\newtheorem{example}[theorem]{Example}

\theoremstyle{remark}
\newtheorem{remark}[theorem]{Remark}

%% ---- Custom Commands --------------------------------------
\newcommand{\R}{\mathbb{R}}
\newcommand{\N}{\mathbb{N}}
\newcommand{\Z}{\mathbb{Z}}
\newcommand{\E}{\mathbb{E}}
\newcommand{\Prob}{\mathbb{P}}
\newcommand{\norm}[1]{\left\lVert #1 \right\rVert}
\newcommand{\abs}[1]{\left\lvert #1 \right\rvert}
\newcommand{\inner}[2]{\left\langle #1,\, #2 \right\rangle}
\newcommand{\ie}{\textit{i.e.}\xspace}
\newcommand{\eg}{\textit{e.g.}\xspace}
\newcommand{\etal}{\textit{et al.}\xspace}

%% ---- Caption Formatting -----------------------------------
\captionsetup{
  font=small,
  labelfont=bf,
  format=hang,
  justification=justified
}

\usepackage{lettrine}

%% ---- Header / Footer --------------------------------------
\usepackage{fancyhdr}
\pagestyle{fancy}
\fancyhf{}
\fancyhead[R]{\small\thepage}
\renewcommand{\headrulewidth}{0.4pt}

%% ============================================================
%%  FRONT MATTER
%% ============================================================

\title{%
	\vspace{-1.5em}%
	\large\textbf{Explorations of the Cauchy Momentum Equation.}%
}
\author{By Amlal El Mahrouss}
\date{\today}

%% ============================================================
\begin{document}
%% ============================================================

\maketitle
\tableofcontents
\newpage

%% ============================================================
\section{Abstract:}
%% ============================================================

\lettrine[lines=1, lhang=0.1, loversize=0.2]{T}HE research paper will cover the convective form of the Cauchy Momentum Equation as stated per Navier-Stokes and its applications in Machine Learning. And to conduct several observations on the said equation.\\
Results and more remarks will be given throughout the paper. Prior knowledge about the Navier-Stokes equations and their derivations is assumed.\\ \\
\textbf{Keywords:} Cauchy Momentum Equation, Navier-Stokes, Machine Learning, Convective Form, Fluid Dynamics, Mathematical Analysis.

%% ============================================================
\section{Introduction and Definitions before Proofs:}
%% ============================================================

\lettrine[lines=1, lhang=0.1, loversize=0.2]{W}E shall first introduce the definitions and remarks in a mathematically sound manner before proceeding towards proofs and graphs.\\Moreover the definitions stated in this preprint will be used later in the paper to provide examples on use cases for this paper. Now going through the definitions:
\subsection{Introduction of the Cauchy Momentum Equation:}
\begin{definition}
There exist by setting the cauchy stress tensor $\sigma$ to be the sum of a viscosity term $\tau$ and a pressure term $-p\, \mathbf{I}$, defined as the volumetric stress, the following equation:
\begin{equation}
\label{eq:cauchy}
  \rho \frac{\mathbf{D}\mu }{\mathbf{D}t} = - \nabla p + \nabla \cdot \tau + \rho \, \textbf{a}.
\end{equation}
\end{definition}

%% ============================================================
\subsection{Remarks and Definitions:}
%% ============================================================

\lettrine[lines=1, lhang=0.1, loversize=0.2]{T}HE following remarks are assumed to be true and will be used in the next sections of this paper. The remarks are as follows:

\begin{remark}
\label{rem:1}
  Let the material derivative be defined as: $\frac{\mathbf{D}\mu}{\mathbf{D}t}$.
\end{remark}
\begin{remark}
\label{rem:2}
  Let the density variable be defined as: $\rho$.
\end{remark}
\begin{remark}
\label{rem:3}
  Let the flow density variable be defined as: $\mu$.
\end{remark}
\begin{remark}
\label{rem:4}
  Let the divergence vector operator, defined as: $\nabla \cdot$.
\end{remark}
\begin{remark}
\label{rem:5}
  Let the pressure variable, defined as: $p$.
\end{remark}
\begin{remark}
\label{rem:6}
  Let the time variable, defined as: $t$.
\end{remark}
\begin{remark}
\label{rem:7}
  Let the deviatoric stress tensor operator, defined as: $\tau$.
\end{remark}
\begin{remark}
  Let the body accelerations defined as $\mathbf{a}$ acting on the continuum.
\end{remark}
\begin{definition}
Assuming we use the defintion per Remark~\ref{rem:7}, redefining Equation~\ref{eq:cauchy}, that yields:
\begin{equation}
  \rho \frac{\mathbf{D}\mu}{\mathbf{D}t} = - \nabla p + \nabla \cdot \tau + \rho \, \textbf{a}.
\end{equation}
Assuming the variable of integration to be Remark~\ref{rem:7} between the open interval $[\mathbb{R}^{+}, +\infty)$:
\begin{equation}
  \label{eq:integral:1}
  \int_{\mathbb{R}^{+}}^{+\infty} \rho \frac{\mathbf{D}\mu}{\mathbf{D}t} \, dt = - \int_{\mathbb{R}^{+}}^{+\infty} \nabla p + \nabla \cdot \tau + \rho \, \textbf{a} \, dt.
\end{equation}
\end{definition}
\begin{theorem}
  \label{ex:1}
  Let the Equation~\ref{eq:integral:1} in $\mathbb{R}$, within an interval of $\mathbb{R}^{+} \land +\infty$:
  \begin{equation}
      - \int_{\mathbb{R}^{+}}^{+\infty} \nabla p + \nabla \cdot \tau + \rho \, \textbf{a} \, dt.
  \end{equation}
  Applying the decomposition theorem for the Equation~\ref{eq:integral:1} for the two intervals $[\mathbb{R}^{+}, x] \land [x, +\infty)$:
  \begin{equation}
      - \int_{\mathbb{R}^{+}}^{x} \nabla p + \nabla \cdot \tau + \rho \, \textbf{a} \, dt \pm \int_{x}^{+\infty} \nabla p + \nabla \cdot \tau + \rho \, \textbf{a} \, dt.
  \end{equation}
\end{theorem}
\begin{definition}
  Let $\mathbb{R}^{+}_0$ be the first element of the set $\mathbb{R}$.
\end{definition}
\begin{definition}
  Let $\mathbb{R}^{+}_p$ be the element of index $p \in \mathbb{N}^{+}$ in the set $\mathbb{R}$.
\end{definition}
\begin{remark}
  There shall be $p \geq 0$, and $\mathbb{R}^{+}_0 \neq \mathbb{R}^{+}_p $ and $ \mathbb{R}^{+}_p > \mathbb{R}^{+}_0$.
\end{remark}
\begin{definition}
  Let the following construction applied from the interval of values of $x$ between $\mathbb{R}^{+}_0$ to $\mathbb{R}^{+}_{p}$, solving for $x=\mathcal{R}$:
  \begin{equation}
    \mathbb{R}^{+}_0 < \mathcal{R} < \mathbb{R}^{+}_{p}.
  \end{equation}
  Reorganizing the inequality gives the following inequality:
  \begin{equation}
    \mathbb{R}^{+}_0 < \mathcal{R} < \mathbb{R}^{+}_{p}.
  \end{equation}
  Imposing a variable $t \in \mathbb{R}$ and a limit of $\mathbb{R}^{+}_0 < \mathcal{R}$ from a value $n \to 0$ would be as the index of $\mathbb{R}^{+}_0$ is zero:
  \begin{equation}
    \lim_{t \to \infty} \Bigl( \mathbb{R}^{+}_n < \mathcal{R} < \mathbb{R}^{+}_{p} \Bigr), \, p \neq 0, \, p > n.
  \end{equation}
  The limit has been invoked per \ref{eq:integral:1} to establish non-degeneracy cases for \ref{eq:cauchy} and \ref{eq:integral:1}. Moreover tigheting the interval of values would turn out to serve us well in the next sections about applied examples of the integral.
\end{definition}

%% ============================================================
\subsection{Final Remarks and Definitions before Proofs:}
%% ============================================================

\begin{definition}
\label{def:plane-1}
Let us define the integral of $[\mathbb{R}^{+}, \infty)$ as:
\begin{center}
\fbox{$-\int_{\mathbb{R}^{+}}^{+\infty} \nabla p + \nabla \cdot \tau + \rho \, \textbf{a} \, dt.$}
\end{center}
For a non-trivial non-linear plane in $\mathbb{R}^{+}$.
\end{definition}
\begin{remark}
\label{rem:plane-1}
The plane per \ref{def:plane-1} is assumed to be non-degenerate and of fourth dimensional space.
\end{remark}

\begin{corollary}
Per \ref{def:plane-1} and \ref{rem:plane-1}, there exist a set of degenerate planes in $\mathbb{R}^{+}$, that are of fourth dimensional space, and are non-trivial and non-linear such that:
\begin{equation}
  \operatorname{Plane}_{n}(x, y, z, t) \notin \mathbb{R}^{+}\, , \forall n \in \mathbb{N}^{+}.
\end{equation}
\end{corollary}

%% ============================================================
\section{Proofs and Applications in a 4D-plane:}
%% ============================================================

%% ============================================================
\subsection{Definitions, Corollaries, And Proofs:}
%% ============================================================

\begin{definition}
\label{def:plane-4d}
A 4D-plane is a 4 manifold in a convective form of the Equation \ref{eq:integral:1} assuming Equation (\ref{eq:integral:1}) to be continuous within the interval: $\mathbb{R}^{+} \land +\infty$.
\end{definition}

\begin{corollary}
\label{cor:plane-4d}
Such definition per Definition \ref{def:plane-4d} would then lead to the assumption that Equation \ref{eq:integral:1} hold for such interval: $\mathbb{R}^{+} \land +\infty$.\\
Thus for Equation \ref{eq:integral:1}, there exists a 4 manifold form of Equation \ref{eq:integral:1}. 
\end{corollary}

\begin{proof}
  Applying a proof by induction per Corollary \ref{cor:plane-4d}, Definition \ref{def:plane-4d}, Theorem \ref{ex:1}, and Equation \ref{eq:integral:1}:\\
  (i) Applying the corolloary \ref{cor:plane-4d} per Definition \ref{def:plane-4d}, there is then a non-degenerate form of Equation \ref{eq:integral:1}.\\
  (ii) Thus per (i), there exist for the Equation \ref{eq:integral:1} a convective form of Equation \ref{eq:integral:1}.\\
  (iii) Per the Theorem \ref{ex:1}, there exist then an identity of Equation \ref{eq:integral:1} in $\mathbb{R}^{+} \land +\infty$.\\
  (iv) Thus completing the proof of the existence of a convective form of Theorem \ref{eq:integral:1} in $\mathbb{R}^{+} \land +\infty$.\\
  (v) We, by rigor, restate the following Equation \ref{eq:integral:1}:
  \begin{equation}
    \fbox{$\int_{\mathbb{R}^{+}}^{+\infty} \rho \frac{\mathbf{D}\mu}{\mathbf{D}t} \, dt = - \int_{\mathbb{R}^{+}}^{+\infty} \nabla p + \nabla \cdot \tau + \rho \, \textbf{a} \, dt.$}
  \end{equation}
\end{proof}

%% ============================================================
\section{Conclusion and Additional Remarks:}
%% ============================================================

Per the stated abstract, the research paper has fulfilled its purpose, additional graphs and results will be published after writing.

\nocite{*}
\bibliographystyle{acm}
%% TODO: More references.
\bibliography{pinter2010book, simons1968, weil1967basic, temam2024navier, munkres}

%% ============================================================
\end{document}
%% ============================================================